\section{Perpetual Motion}

\begin{quotex}
The present age is a critical one and interesting to live in. The civilisation characteristic of Christendom has not disappeared, yet another civilisation has begun to take its place. The shell of Christendom is broken. The unconquerable mind of the East, the pagan past, the industrial socialistic future confront it with their equal authority. Our whole life and mind is saturated with the slow upward filtration of a new spirit---that of an emancipated, atheistic, international democracy.
\flright{\textsc{George Santayana}, \emph{The Winds of Doctrine (1913)}}
\end{quotex}


\paragraph{Motion Detector}


While discussing the discoveries of the cognitive psychologist James Gibson, Wolfgang Smith pointed that we don't perceive motion as a series of independent frames, but rather we grasp it ``all at once''. Music, too, is experienced ``all at once''; that is, songs, concertos, symphonies, etc., are recognized as a ``whole'', not as a series of unrelated notes.

This experience is so common, it is barely noticed, if at all. It is the cause of the ``gambler's fallacy''. A gambler, having seen Red come up 10 consecutive times on the roulette wheel, believes that a Black is ``due'' to come up. The gambler is deluded into thinking there must be a relationship between the rolls of the roulette wheel.

\paragraph{Civilisational Decline}


When it comes to the larger question of the course of human events, the opposite is the case. Then only individual events are noticed, but few are able to see these events in their Wholeness as part of a larger pattern lasting, perhaps, generations. Early last century some astute minds were able to see the ``motion'' and ``direction'' of events. I am thinking of Oswald Spengler, for example, but also Rene Guenon and Julius Evola. Writing in 1913, George Santayana, who famously declared that ``Those ignorant of history are doomed to repeat it,'' had his own astute observations of events. Unfortunately, in our day, it seems that we are doomed to repeat Santayana's phrase ad infinitum, while still remaining ignorant of history. The civilization in decline, more complete now than then, is called \textbf{Christendom}. Santayana identifies the forces that have been, and still are, replacing it:

\begin{itemize}


\item The mind of the East

\item Paganism

\item Atheism

\item Socialism

\item Globalism
\end{itemize}


This list includes those movements whose goal is to ``save the West'', as well as those explicitly committed to its continued destruction. Here is the first distinction.

\begin{itemize}


\item Those committed to destruction understand they are destroying Christendom.

\item Those allegedly committed to preservation don't understand that they need to preserve Christendom.
\end{itemize}


\subparagraph{Turn to the East}


Not much needs to be said about the turn to Eastern religions. Thanks to the researches of Rene Guenon, we realize how pointless and unnecessary that is. There is nothing in an Eastern religions that cannot be found in the West; if you don't recognize it in the West, then you won't recognize it in the East. I include Eastern Orthodoxy in the group. I have great respect for Orthodoxy as it is appropriate to those born into it. But once again, the same principle applies: esoterically, there is no difference between Latin and Greek Christians. Perhaps your complaint is that the Roman church has ``changed''. Well, if someone has taken something that belongs to you, then a man fights to get it back.

\subparagraph{Paganism}


Regarding the neopagans, I am tempted to say, ``forgive them, for they don't know what they do,'' but the thought quickly recedes. The province of yobs and the semi-intelligent, they yearn for a past the preceded Christendom. Santayana identifies them as destroyers, so they are doing the enemy's work for them.

\subparagraph{Atheism, Socialism, Globalism}


This group, which we can't dwell on at this point, explicitly intends to overturn civilization. They have the upper hand.

\paragraph{Common Ground}


C S Lewis, in an essay titled\footnote{\url{https://theimaginativeconservative.org/2015/12/c-s-lewis-on-modern-man-and-his-categories-of-thought.html}} \textit{Modern Man and His Categories of Thought} (1946), identified the modes of thought that separates the modern and traditional minds.

Speaking of Jews, Gentiles, and Pagans, he wrote: ``all three classes believed in the supernatural... All were conscious of sin and feared divine judgement... All three believed that the world had once been better than it now was.'' The Pagans, in particular, had ``reverence for heroes, ancestors, and ancient lawgivers.''

That is why, in the Middle Ages, the Christian Scholastics, the Jew Maimonides, or the Muslim Avicenna, could all draw on the thought of the Pagans Plato and Aristotle; thus, they could have mutually intelligible conversations with each other, despite their exoteric differences. If you listen in on their discussions and then compare them to the one between Bishop Barron and Ben Shapiro, you cannot help but be struck by the decline in the quality of thought.

\paragraph{The Devil Will Find You}


In a podcast on \textit{In Search of Chartres}\footnote{\url{https://isoc.ws}}, Peter Chojnowski went into a rant against \textbf{Wolfgang Smith}. He assumed that Dr. Smith is messing with Ouija boards, Tarot spreads, and the ``intermediary realm''. Dr. Chojnowski claims that the realm is inhabited by devils. I am sure there are some aspects of Dr. Smith to be critical about, but Peter is really not capable of doing it well. I will leave him with one warning, and this applies to whoever is reading this: \emph{The Devil will find you in whatever ``realm'' you think you are hiding in}. Keep your guard up.

\paragraph{Do What you Will}


\begin{quotex}
\emph{Der Mensch kann zwar tun, was er will, aber er kann nicht wollen, was er will.}

A man can certainly do what he wills, but he can't will what he wills.
\flright{\textsc{Arthur Schopenhauer}}
\end{quotex}


If you are wondering how you can save civilization, the Saturnine philosophers are not too sanguine about your prospects. Santayana is the worst of all, believing that humans are just the playthings of material forces. Schopenhauer concedes that you can choose between a lager or an IPA, but you can't choose not to enjoy a beer. Other than that, a blind cruel will is behind everything. Contra him, we do see pockets of order here and there.

Spinoza claims that we simply are unaware of the forces that guide out will; salvation is becoming aware of those forces. This is much more than simplistic notions of self-knowledge. For example, a man nay ``like'' a politician and get excited about his policies. So the man thinks he knows himself. But what he does not understand is why he likes that politician. That is real self-knowledge.

\paragraph{Poker Training}


In my ideal curriculum, everyone would take a semester of poker --- using your own money, of course; ideally, your tuition money to make it spicy. Then you will gain self-knowledge as well as knowledge of others. The worst players think they will hit the flush or inside straight. Although they seldom do, the thrill from doing it once is too much to resist. But, as the song says, you need to know when to hold them and when to hold them.

The same principle applies to an online argument. You need to know when you've lost an argument and give it up. The yobs always try to have the last word. There needs to be a fee to post a comment; skin in the game changes everything.

\paragraph{Infinity}


Every once in a while, I awaken in the middle of the night and contemplate infinity. Unlike a Nietzsche who believes in the infinite return of the same, I understand that the infinite has too much variety to be restricted to like that. Besides, I have a large task ahead of me that may require infinity.

You may be pleased to know that every infinite set has a subset equivalent to $\omega$. That is, the natural numbers will belong to every possible world.

Moreover, every set can be well ordered. So just when you suppose that the chaos cannot get any worse, you are not seeing the world the way that God does.


\flrightit{Posted on 2021-05-02 by Cologero}

\begin{center}* * *\end{center}

\begin{footnotesize}
\begin{sffamily}
\texttt{N.W. Flitcraft on 2021-05-02 at 22:57 said:}

I posted this on Chojnowski's blog, though I don't think it's been approved yet:

From the fact that the intermediary domain---what the Orthodox call the ``aerial world''---is where the exorcist communicates with demonic entities, it does not follow that all one can say about this domain is that demons operate there. That would be like saying that because ICE spends most of their time in Texas and Arizona, all that one will find in Texas and Arizona are illegal immigrants.

\hfill

\texttt{Auld Wat on 2021-05-03 at 16:36 said:}

Nietzsche's concept of Eternal Return is useful as a stepping stone to greater things. Likewise Herman Hesse's Steppenwolf sees the infinite possibility of his personalities in the fragmented mirror. It is a recognition of a higher state of being above the monotony of eternity and the disorder of unconsciousness.

\hfill

\texttt{Jorgen B on 2021-05-23 at 23:29 said:}

``There is nothing in an Eastern religions that cannot be found in the West''

Anything Buddhism can do NeoPlatonism can do, and since it accepts God, it can be merged with Christianity. That may be called ``heresy'' but its no worse than any other commonly accepted heresy; since NeoPlatonic Christianity will of necessity lead to amillenialism its much better than the dispensationalist heresy.

\hfill

\texttt{Dimitri on 2022-05-08 at 00:44 said:}

Genuine question for any at Gornahoor that wish to respond: since we acknowledge that the so-called pagans of the past are just as integral to Western Tradition as Christianity is, and pagan religion and Christianity had become synthesized, does it follow that it is acceptable for a Christian to pray to Jove as he would to Michael, or have devotions to Ceres as they would to Mary? If these are all divinities and are all part of Western Tradition, and we are to love the pagan past, this should not be problematic, except perhaps on the exoteric level, no? What would the ``Hermetic'' response to this be?

\hfill

\texttt{Cologero on 2022-05-08 at 09:43 said:}

That is not a genuine question, Dimitri, as I doubt he really gives you agita. Perhaps you could look up the word ``past'' in the dictionary, as that might help.

\hfill

\texttt{Dimitri on 2022-05-08 at 12:45 said:}

I confess I did have to look up the meaning of ``agita'' to understand what you meant. I wouldn't know whether Olympians can provide that, as I don't pray to them myself (maybe the Romans would have disagreed?), but I do want to figure out what role they would have for us now. I'd like to think they have a place somewhere rather than just being a thing of the past. Perhaps they go by different names these days. Do forgive any perceived ignorance on my part, by the way.

\end{sffamily}
\end{footnotesize}

